\chapter{Las fincas y sus nombres}

\label{cap:fincas}
\begin{cautela}
Antes de empezar conviene decir qué hay debajo de todo lo que sigue.

Este capítulo reconstruye quién tuvo la tierra de este departamento desde 1908, con los
nombres que el archivo entrega: fincas, deslindes, remates, sucesiones. \textbf{Esa no es la
capa más antigua.}

En 2011 el Instituto Nacional de Asuntos Indígenas relevó la ocupación tradicional de la
\textbf{Comunidad Originaria Kolla Kondorwaira} y estableció que la zona alta de la cuenca
del río La Caldera es territorio de ocupación actual, tradicional y pública de esa comunidad.
El relevamiento fue aprobado por resolución en 2013 y ratificado por dos organismos
provinciales en 2023. El capítulo \ref{cap:resistencias} lo documenta.

De modo que \textbf{los partidos de 1909, las fincas del deslinde de 1932 y las cuarenta y
siete partidas del revalúo de 1944 son una capa que se superpone a otra anterior}, que este
libro no puede reconstruir porque no dejó el tipo de rastro que el Boletín Oficial publica.

Y hay una consecuencia práctica que conviene tener presente en todo lo que sigue: cuando este
capítulo dice que una finca ``se deslindó'' o que un partido ``se creó'', describe actos del
Estado provincial sobre un territorio que ya estaba ocupado. \textbf{El archivo registra la
capa que él mismo produjo.}
\end{cautela}

\section{Dos topónimos que son títulos}

Vaqueros toma su nombre de Francisco Vaquero, poseedor de las tierras donde hoy se
asienta la localidad, que hacia 1690 desarrolló allí agricultura, ganadería y molienda de
cereales. El sitio funcionó como parador del comercio de ganado mular hacia los pueblos
del Pacífico y llegó a ser, según la memoria institucional del municipio, el polo de
molienda más importante de los alrededores de la ciudad de Salta. En su museo popular se
conserva un molino hidráulico de piedra cuyo funcionamiento está registrado en 1735.

El topónimo del municipio es, entonces, el apellido de un propietario colonial. La unidad
político-administrativa que hoy aprueba loteos lleva inscripta en su nombre la estructura
de tenencia que la precedió en tres siglos.

Sobre el origen del nombre de la cabecera departamental circulan dos versiones —un
misionero jesuita de apellido Caldera, o una antigua fundición de plomo— y ninguna es
concluyente. Lo estructuralmente relevante es la posición: La Caldera está sobre el camino
real a Jujuy, hoy Ruta Nacional 9, donde el Valle de Lerma se cierra y comienza la
cornisa. Conviene sí dejar anotado, porque el capítulo \ref{cap:siglo} lo documenta, que en el
departamento funcionó efectivamente una fundición metalúrgica de plomo, la de Mojotoro,
al menos entre 1917 y 1925. Es demasiado tardía para explicar un topónimo colonial, pero
prueba que la actividad no era ajena al lugar.

\subsection*{Las fincas que la norma enumera}

El artículo 6º del proyecto de ordenanza del Plan de Desarrollo Urbano Ambiental enumera
las principales localidades y parajes que integran el municipio de La Caldera, y esa
enumeración es el mapa de tenencia histórica consignado en el instrumento de planeamiento de
2015: La
Calderilla, Finca El Gallinato, Dique Campo Alegre, Finca San Alejo, Finca Los Porongos,
Finca San Ignacio, Finca Las Cañas, Potrero de Castilla, Garabatal y Yacones.

Cinco fincas nombradas como unidades territoriales integrantes del municipio, más un
potrero. El mismo artículo consigna que el municipio abarca 1.200 km², equivalentes al 82
por ciento de la superficie departamental y al 10 por ciento de la provincial. Ninguna de las dos proporciones resiste el cotejo, y el motivo es que \textbf{el departamento no tiene una superficie única}: circulan al menos seis cifras, y cada una sale de una fuente distinta.

\begin{ficha}{La superficie del departamento La Caldera según la fuente}
\begin{center}\small
\begin{tabular}{@{}r p{5.2cm} p{5.2cm}@{}}
\toprule
\textbf{km²} & \textbf{Fuente} & \textbf{Qué es} \\
\midrule
867 & Anuario Estadístico de la Provincia, cuadros censales 2001--2010 & Cifra estadística anterior, heredada del INDEC \\
1.076 & Radios del Censo 2022 (elaboración propia, capítulo \ref{cap:opacidad}) & Suma de los veintiocho polígonos censales \\
\textbf{1.091,9} & \textbf{INDEC, Censo 2022, cuadro 2.17}, que la toma de la determinación oficial de superficies del Instituto Geográfico Nacional (tabla 29) & \textbf{Cifra de referencia de este libro}; la misma que publica el Senado provincial redondeada a 1.092 \\
1.092,1 & Capa «Departamento» de las Capas SIG del Instituto Geográfico Nacional, con fuente declarada en la Dirección General de Inmuebles de Salta, descargada en septiembre de 2026 (medición propia en proyección equivalente) & Polígono cartográfico vigente; la suma de las capas de municipios del mismo organismo ---772,3 km² para La Caldera y 319,8 para Vaqueros--- da la misma cifra. Que la fuente declarada sea Inmuebles no lo convierte en el límite legal \\
1.109 & Capa de límites del Instituto Geográfico Nacional, 1:250.000 (medición propia, capítulo \ref{cap:opacidad}) & Polígono cartográfico medido en proyección; difiere de la cifra oficial del propio IGN por el método \\
1.440 & Sitio de la Municipalidad; fundamentos de un proyecto de ley de 2003 publicado en el Boletín Oficial & Sin fuente técnica declarada; suma 1.200 de La Caldera y 240 de Vaqueros \\
\bottomrule
\end{tabular}
\end{center}
Una versión anterior del capítulo \ref{cap:tierra} usaba además unas 146.300 hectáreas como superficie departamental para calcular la cobertura de las explotaciones agropecuarias; es la mayor de todas y no coincide con ninguna fuente cartográfica.
\end{ficha}

\textbf{Este libro adopta como referencia la cifra con que el INDEC calcula la densidad del Censo 2022: 1.091,9 km², unas 109.200 hectáreas}, que da 11,3 habitantes por kilómetro cuadrado. Las demás se citan con su fuente cuando se usan. Las tres mediciones cartográficas propias ---1.076, 1.092,1 y 1.109--- quedan a menos del dos por ciento de la de referencia, lo que la respalda; la capa vigente del IGN coincide con ella a 0,2 km², pero esa coincidencia puede deberse a que ambas cifras salen del mismo polígono y no prueba por sí sola que sea el límite correcto; las de 867 y 1.440 quedan lejos, por defecto y por exceso, y la de 146.300 hectáreas, más lejos todavía.

\begin{cautela}
La cifra está confirmada en su fuente primaria: el cuadro 2.17 de los resultados definitivos del Censo 2022 asigna al departamento 1.091,9 km² y a la provincia 155.340,5, y aclara en nota que los valores provienen de la «Determinación de la superficie correspondiente al Territorio Continental, Antártico e Insular de la República Argentina» del IGN. Es una superficie estadística y cartográfica, no la legal. \pendiente{constancia de la Dirección General de Inmuebles sobre la superficie legal del departamento y de cada municipio}
Con esa referencia, la primera proporción del plan se vuelve imposible: un municipio de 1.200 km² no cabe en un departamento de 1.091,9. La segunda ya lo era: el 10 por ciento de Salta serían unos 15.500 km². El departamento representa cerca del 0,7 por ciento de la provincia. La superficie legal ---la que fija el límite trazado por las Leyes 2131, 4365 y 4987--- es la de la Dirección General de Inmuebles, y es la que corresponde pedir.
\end{cautela}

\section{Los partidos de 1909}

El 12 de febrero de 1909 el Poder Ejecutivo designó comisarios auxiliares para cada
partido del departamento. El decreto es, sin proponérselo, el mapa administrativo completo
del territorio un siglo antes de los loteos.

\needspace{6\baselineskip}
{\sffamily\bfseries\footnotesize Partidos del departamento La Caldera y sus comisarios auxiliares, 1909\par}
\vspace{0.5em}
{\small
\setlength{\LTpre}{0.6em}\setlength{\LTpost}{1.2em}
\begin{longtable}{@{}ll@{}}
\toprule
\textbf{Partido} & \textbf{Comisario auxiliar} \\
\midrule
\endhead
\textbf{San Alejo} y Santa Rufina & Ángel Fernández\index{Fernández, Ángel} \\
La Calderilla & Genaro Molina\index{Molina, Genaro} \\
\textbf{Vaqueros} & Adolfo Agote \\
Los Sauces y Angostos de Arias & Julio Mercado \\
La Despensa & Francisco Leeds \\
\textbf{Los Porongos} & Teodoro Reyes \\
Peñones y Chalchanio & Liborio Guerrero\index{Guerrero, Liborio} \\
Mojotoro & Modesto Mantilla \\
\textbf{Potrero de Castillo} & Zacarías Maidana \\
Potrero de Valencia & Abraham Fernández \\
\textbf{Potrero de Gallinato} & Silverio Cáceres \\
Lesser & Marcelo Lamas \\
Los Yacones & Victoriano Vivero \\
\bottomrule
\end{longtable}
}

\textit{Comisario suplente de policía del departamento: Jacinto Guerrero\index{Guerrero, Jacinto}.}

El comisario titular no está en ese decreto sino en uno anterior. El 23 de diciembre de 1908
el Poder Ejecutivo nombró los comisarios de policía de los departamentos de campaña para el año
siguiente, y el primero de la lista es «Caldera, á don Escolástico Aparicio»\index{Aparicio, Escolástico M.} (B.O.\ Nº 22,
30/12/1908, p.\ 105). El artículo 2º le encargaba recibir de su antecesor el archivo y los enseres
«bajo de inventario» y \textbf{proponer, por intermedio del jefe de policía, a las personas que debían
desempeñar las comisarías de partido}. Los trece comisarios auxiliares de febrero son, por ese
procedimiento, los que el comisario del departamento debía proponer; el decreto de febrero no dice
quién los propuso.

Seis semanas antes, el 12 de noviembre, un decreto del Ministerio de Hacienda había nombrado a
«Escolástico M.\ Aparicio» \textbf{clasificador de patentes generales} de Caldera para 1909 (B.O.\
Nº 10, 18/11/1908, p.\ 54): el que clasificaba a los contribuyentes del departamento para el
impuesto de patentes. Los reclamos contra esa clasificación los resolvía un jurado del mismo
departamento, nombrado el mismo día: Adolfo Vidal, Abelardo Figueroa y Carlos Murúa, con
Augusto Regis de suplente (apéndice \ref{cap:personas}). \textbf{En el primer año del Boletín, el círculo de personas que el
Estado provincial designa en La Caldera es chico}: el que clasificaba el impuesto pasó a mandar
la policía, y dos de los nombrados para el jurado de reclamos, Figueroa y Regis, reaparecen en la
comisión municipal de 1909 y en la comisión de defensas de 1910.

\subsection*{La otra división del mismo año}

Dos meses después, el mismo territorio se dividió de otra manera y para otra cosa. El 8 de abril
de 1909, en «\textit{este pueblo de Caldera, capital del departamento del mismo nombre}», el juez
de paz Carlos Matienzo\index{Matienzo, Carlos} y dos vecinos propietarios ---Daniel
Linares\index{Linares, Daniel} y Jacinto Guerrero\index{Guerrero, Jacinto}--- nombraron los
\textbf{veedores} del artículo 33 del Código Rural, y la lista se publicó por orden del
Ministerio de Gobierno (B.O.\ Nº 51, 1909, h.\ 2).

\needspace{6\baselineskip}
{\sffamily\bfseries\footnotesize Secciones y veedores del Código Rural, 1909\par}
\vspace{0.5em}
{\small
\setlength{\LTpre}{0.6em}\setlength{\LTpost}{1.2em}
\begin{longtable}{@{}ll@{}}
\toprule
\textbf{Sección} & \textbf{Veedor} \\
\midrule
\endhead
San Alejo y Santa Rufina & Andrés V.\ Matienzo \\
La Despensa & Abdón Choque \\
Los Porongos & Teodoro Reyes \\
Los Peñones & Livorio Guerrero\index{Guerrero, Liborio} \\
Chalcanio & Héctor Echenique \\
\textbf{Vaqueros} & Higinio Juárez \\
\textbf{Wierna} & Darío Moya \\
\textbf{Potrero de Castillo} & Laureano Quipildor \\
Gallinatos & Silverio Cáceres \\
Valencia & Benancio Ruiloba \\
Lesser & David Segovia \\
Los Yacones & Antonio Vera \\
Los Sauces, Angosto de Arrieta y Quecería & Miguel Ceballos \\
Bandera Angosta y Mojotoro & Daniel Arámbulo \\
Mojotorillo y Chaguadero & Desiderio Gumiel \\
\bottomrule
\end{longtable}
}

\textbf{Las dos listas no coinciden, y la diferencia importa.} La policial tiene trece partidos;
la rural, quince secciones. La rural agrega \textbf{Wierna}, \textbf{Bandera Angosta},
\textbf{Mojotorillo}, \textbf{Chaguadero} y \textbf{Quecería}, que la policial no nombra, y
fusiona Los Sauces con el Angosto de Arrieta. El mismo territorio, el mismo año y dos mapas
distintos según para qué se lo dividiera: vigilar personas o vigilar ganado. Es la forma más
antigua de algo que este libro documenta un siglo después ---el departamento tiene tantos
límites como organismos lo dividan--- y la primera vez que \textbf{Wierna} figura como unidad
territorial, doce años antes del circuito electoral de 1921.

Un nombre de esa lista vuelve más adelante: el veedor del Potrero de Castillo es \textbf{Laureano
Quipildor}, y en 1925 un pedimento de cateo declara terrenos linderos de «E.\ Quipildor» junto a
la finca Las Nieves, en la misma zona (capítulo \ref{cap:siglo}). Es coincidencia de apellido, no
identidad probada.

Trece partidos. Leída contra el artículo 6º del Plan, la correspondencia es casi completa:
El Gallinato, San Alejo, Los Porongos y Potrero de Castilla —cuatro de las seis unidades
que el proyecto de ordenanza de 2015 enumera— son partidos policiales de 1909. Y lo mismo vale para los
demás nombres que este libro persigue. La Finca Potrero de Gallinato, sobre la que la
Provincia declaró en 2015 una prescripción adquisitiva para regularizar una escuela, y
donde ese mismo año se aprobó el club de campo Chacras del Gallinato, es el partido de
Silverio Cáceres. El río Porongal, por cuyo talweg se delimitaba un proyecto de reserva
natural en 2023, es el partido de Teodoro Reyes. El partido de Potrero de Castillo es donde
en diciembre de 1908 se pidió cateo de cobre. Lesser, Los Yacones, Mojotoro, La Calderilla
y Vaqueros —hoy parajes, cursos de agua, canteras o un municipio entero— tenían comisario
propio.

\textbf{Con Mojotoro, sin embargo, las fuentes del propio 1909 no coinciden.} El decreto de abril
que nombra al comisario de la estación lo hace «\textit{comprensión del departamento de la
Caldera}», y una causa criminal del mismo año describe a un procesado como domiciliado en
«\textit{Mojotoro, jurisdicción del departamento de Campo Santo}». El límite pasaba por ahí y los
dos organismos lo ubicaban de distinto lado.

Dos años después, en abril de 1912, el decreto que renueva los comisarios auxiliares hace
explícito el mecanismo que el de 1909 sólo mandaba: nombra a los de doce partidos del
departamento «\textit{de acuerdo cón las propuestas presentadas por los Comisarios de Policia}»
(B.O.\ Nº 330, 1912, h.\ 2). Los nombres vuelven a ser los mismos ---Ángel
Fernández\index{Fernández, Ángel} en San Alejo, Juan Chuchuy\index{Chuchuy, Juan} en los
Yacones, Teodoro Reyes en los Peñones, Francisco Sosaya en el Angosto de Arrieta, Héctor
Echenique en Chalchañio, Eustaquio Murúa\index{Murúa, Eustaquio} en La Despensa--- y el suplente
del departamento es otra vez Jacinto Guerrero\index{Guerrero, Jacinto}. El mismo decreto trae
una errata que conviene conservar porque muestra cómo se imprimían estos nombres: escribe
«\textit{el Augosto de Arias}» \textit{[sic]} tres renglones antes de «\textit{el Angosto de
Arrieta}».

En abril de 1913 la lista vuelve a tener \textbf{trece partidos} y, por primera vez, nombra
juntos a \textbf{Lesser}, \textbf{Gallinatos} y \textbf{Angosto de Arias} como partidos propios
(B.O.\ Nº 404, 1913, h.\ 2). El decreto se imprime con tres erratas seguidas ---«los
\textit{cindadanos}» \textit{[sic]}, las «\textit{cocisarias}» \textit{[sic]}
«\textit{auxilaiares}» \textit{[sic]}--- y una de sus designaciones dura siete días: el 18 de
abril otro decreto deja sin efecto la de Los Peñones y le da el cargo a \textbf{Enrique
Leconte}\index{Leconte, Enrique}, que era comisario auxiliar de Los Porongos y que, cuatro meses
antes, había pedido el deslinde de la finca Los Porongos o San Luis (apéndice
\ref{cap:dominio}). \textbf{El que mandaba a medir la finca mandaba también la policía del
partido donde estaba.}

\textbf{Tres años después el pleito llega a una sentencia, y tampoco ahí se cierra.} En junio de
1912, en el juicio de los señores Paulucci y Ricchelli (B.O.\ Nº 344, 1912, h.\ 2), el juzgado afirma que la finca Mojotoro
es «\textit{jurisdicción del Departamento de La Caldera (no Campo Santo)}» ---el paréntesis es
del original---, mientras el demandado sostiene lo contrario. En la misma sentencia consta que
la carnicería y el despacho de bebidas del partido de Mojotoro pagaban patente \textbf{a la
Municipalidad de la Caldera}, que llegó a exonerarlos por tres meses de 1909 (B.O.\ Nº 346, 1912, h.\ 1). \textbf{El
municipio cobraba ahí; el límite, sin embargo, se discutía en juicio.} No se concilia: se
registra.

\textbf{Al año siguiente la lista ya había cambiado, y cambió menos de lo que parece.} El
decreto del 2 de junio de 1910 (B.O.\ Nº 161, 1910, h.\ 3 y 4) nombra los comisarios de partido «\textit{de acuerdo con las
propuestas presentadas por el señor comisario de policía del departamento}» y enumera
\textbf{doce}: La Calderilla (José Miguel Rauch\index{Rauch, José Miguel}), San Alejo y Santa Rufina (Ángel
Fernández\index{Fernández, Ángel}), Vaqueros (Jesús Ramírez), Yacones (Victorino Viveros),
Potrero de Castillo (Zacarías Maidana), \textbf{Bandera Angosta y estación Mojotoro} (Antonio
Maglioni), Potrero de Gallinato (Silverio Cáceres), Potrero de Valencia (Abraham Hernández),
Angostura de Arias de los Sauces (Virgilio Arévalo), Porongos (Teodoro Reyes), Peñones y
Chalchañio (Liborio Guerrero\index{Guerrero, Liborio}) y \textbf{Angostura de Arrieta y Campo
Alegre} (Francisco Sosaya). En julio, un segundo decreto completa los dos que faltaban, Lesser
y La Despensa (B.O.\ Nº 171, 1910, h.\ 3). \textbf{Las dos novedades que este libro atribuía a 1911 ---la estación Mojotoro
y el partido de Campo Alegre--- ya están en 1910}, y el decreto del 23 de marzo de 1911, que
vuelve a nombrar trece partidos, las confirma en lugar de estrenarlas. Que Mojotoro pase de
partido a estación indica que el ferrocarril ya organizaba el territorio por su cuenta.

Y hay un detalle que anticipa en veintisiete años algo que el capítulo documenta más adelante.
El comisario auxiliar de La Calderilla nombrado en junio de 1910, \textbf{José Miguel Rauch}\index{Rauch, José Miguel},
aparece en octubre del mismo año como \textbf{damnificado en una causa por hurto de ganado
sustraído de su finca en La Calderilla} (B.O.\ Nº 200, 1910, h.\ 1): el comisario del partido era propietario en el partido
que vigilaba.

Sesenta y un años después, la Ley 4486 declararía de utilidad pública las tierras
necesarias para construir allí el embalse que abastece a la ciudad de Salta. El topónimo,
entonces, es anterior a la obra: la obra pública heredó los nombres del territorio que
intervino.

Corresponden dos precauciones. Un partido policial no es una finca en sentido registral:
son categorías distintas —una de jurisdicción, otra de dominio— y que compartan nombre
indica que la finca dio nombre al partido, o al revés, pero no que sus límites coincidan.
Y los nombres tienen variantes: «Potrero de Castillo» en 1909 y «Potrero de Castilla» en la
norma de 2015; «Los Porongos» como partido y «Porongal» como río. Son diferencias menores,
atribuibles a la transcripción o al uso, pero conviene registrarlas antes de tratar los
topónimos como idénticos.

La lámina \ref{fig:fincas} lleva esa comparación al catastro vigente. Las parcelas que conservan
el nombre de una finca con actos de dominio publicados entre 1909 y 1945
(apéndice \ref{cap:dominio}) son 294 y suman unas 48.100 hectáreas: \textbf{casi la mitad de la
superficie catastrada del departamento sigue llevando, en el registro de hoy, el nombre que tenía
en el Boletín de hace un siglo}. La coincidencia es de nombre; si los límites son los mismos, sólo
lo dirían los antecedentes dominiales.

\begin{figure}[p]
\centering
\includegraphics[width=\textwidth,height=0.84\textheight,keepaspectratio]{fig-fincas-historicas.png}
\caption[Las fincas del archivo en el catastro vigente]{\textbf{Las fincas del archivo en el catastro vigente.} A: parcelas del departamento cuyo campo «finca» del catastro parcelario coincide con una finca objeto de un acto de dominio publicado entre 1909 y 1945, coloreadas según el año del primer acto (apéndice \ref{cap:dominio}); en gris, las parcelas urbanas. B: el corredor La Caldera--Vaqueros. La coincidencia es de nombre, no de deslinde: el catastro no dice si los límites actuales son los de entonces. Elaboración propia con el catastro parcelario de la Provincia (IDESA, consultado en septiembre de 2026) y la cartografía de radios del Censo 2022 (INDEC). Elaboración propia; dibujo de \texttt{scripts/mapa\_fincas.py} del repositorio \texttt{mapas\_caldera} (apéndice \ref{cap:fuentes}).}\label{fig:fincas}
\end{figure}

\section{La finca Vaqueros, rematada en 1935}

El topónimo del municipio es el apellido de un propietario colonial. La finca que lo lleva
aparece en el archivo en 1935, y aparece en un remate.

\begin{ficha}{Aviso judicial --- B.O.\ Nº 1.582, 3 de mayo de 1935}
Por orden del Juez de Comercio, el martillero Antonio Forcada venderá el 27 de mayo la finca
\textbf{«Vaqueros o Entre Ríos»}, ubicada en La Caldera, «limitando: Norte, Anacleto Toranzo
y Teodolinda de Toranzo; Sud, Francisco López Cross y Mercedes Carrillo; \textbf{Este,
Francisco Urquiza\index{Urquiza, Francisco S.}}; Oeste, \textbf{camino nacional}».

Embargada en la ejecución hipotecaria \textbf{Banco Provincial de Salta contra Ramón
Bracheri}. Base: \$12.000. Seña: 20 por ciento.

\textit{Los apellidos se transcriben literales del escaneo; el del lindero norte admite
leerse Toranzo, Toranzos o Toranzoz, y conviene cotejarlo contra el original.}
\end{ficha}

Cuatro datos.

\textbf{La finca se llama «Vaqueros o Entre Ríos».} El segundo nombre la describe: está
entre ríos, que es exactamente la condición que ochenta años después hará que el río Vaqueros
sea el límite entre dos jurisdicciones y que ninguna administre el conjunto del cauce.

\textbf{Su lindero al este es Francisco Urquiza\index{Urquiza, Francisco S.}}, el mismo que en 1921 hizo deslindar Cerro
de Buena Vista, en 1924 obtuvo la concesión del río Wierna, en 1928 renunció a la Comisión
Municipal y en 1931 le ganó el juicio a la Provincia. En 1935 su propiedad limita con la
finca que da nombre al municipio.

\textbf{Su lindero al oeste es el camino nacional.} La finca da al camino real a Jujuy, hoy
Ruta Nacional 9: el mismo corredor que estructura el mercado de suelo del departamento
noventa años después.

\textbf{Y se remata por una ejecución hipotecaria del Banco Provincial.} Como las dos casas
quinta de 1927 y 1929 y como la estancia Potrero de Gallinato en 1922 y 1930, la tierra de
este departamento cambia de manos, en el registro público, principalmente por deuda. En todo el período relevado hasta esa fecha ---de 1908 a 1935--- este libro encontró un solo aviso de
venta voluntaria de finca, y también fue un remate: el que Alberto, Abel y Adolfo Murúa
pusieron en marcha por su cuenta en octubre de 1933 sobre La Despensa (apéndice
\ref{cap:dominio}). Todo lo demás son remates judiciales.

Es una observación sobre las fuentes antes que sobre el mercado. Las compraventas
particulares no se publican y los remates judiciales sí, de modo que el Boletín muestra
sólo la parte forzada de las transferencias. Pero esa parte forzada es la única que dejó
rastro público, y es sobre ella que se puede reconstruir algo.

\section{Cómo se describía un límite}

En marzo de 1910 se publicó el edicto de mensura de la finca Los Porongos, solicitada por
Francisco Alemán con poder de Mónica Diez de Costas. Sus límites eran, al poniente, «las
lomas altas que lindan con tierras de La Caldera y Sauces, hasta donde se hallan unos
mojones antiguos»; al norte, «la loma que viene dividiendo las tierras de Perico, que hoy
poseen los herederos de don Domingo Iriarte»; al naciente, «siguiendo dicha cuchilla al
encontrarse la loma conocida con el nombre de Las Víboras».

\textbf{Seis años después, la misma finca sale a remate y el aviso sí da la superficie.} En abril
de 1916 el martillero Enrique Sylvester anuncia la venta de «\textit{la mejor y más productiva finca
Los Porongos hoy San Luis}», por la ejecución de Enrique Werner contra \textbf{Humberto Graf von
Garnier Turawa}, con base de \$33.333,33 ---«\textit{o sean las dos terceras partes de su tasación
fiscal}»--- y remate fijado para el 15 de junio en el local del Jockey Bar. Declara
\textbf{5.359 hectáreas, 77 áreas y 83 centiáreas}, según el deslinde de los agrimensores Pfister y
Martin aprobado por el Departamento Topográfico, y repite los mismos límites de 1910 palabra por
palabra, agregando al sud «\textit{la cuchilla que sale del morro alto por el Poniente, conocida por
la del Tabacal o Pelada Muerta, que es la que divide estas tierras de las de la Despensa}» (B.O.\
Nº 600, h.\ 3).

\textbf{Con eso, Los Porongos es la segunda finca del departamento con algo parecido a una cadena.}
En 1910 la hace deslindar Mónica Díez de Costas; el edicto de deslinde de Loma Bola de febrero de
1916 la describe como «\textit{la que fue de Mónica Díaz de Cortes, hoy del conde Garnier de
Turava}»; en abril de 1916 se remata por ejecución contra ese mismo conde; y en 1932 el deslinde de
las cinco fincas de Silvano Ignacio Murúa\index{Murúa, Silvano Ignacio} la registra como «\textit{la
finca y la estancia Los Porongos, de don Daniel Linares}»\index{Linares, Daniel}. \textbf{Y el ejecutante no es un acreedor de paso.} En septiembre de 1912 el deslinde de esa misma
finca ---publicado entonces como «Los Porongos» o «San Luis»--- lo habían solicitado
\textbf{Pablo Enrique Wernner y Enrique Leconte}, y en 1913 un Pablo Enrique Wenner hace deslindar un
terreno propio en el departamento (apéndice \ref{cap:dominio}). Leconte era, ese mismo año, comisario
auxiliar del partido de Los Porongos.

\begin{cautela}
\textbf{Los tres avisos imprimen tres grafías ---Wernner, Wenner, Werner--- y este libro no las
unifica.} Lo que se afirma es lo que está impreso: quien en 1912 pidió el deslinde de esa finca lleva
el mismo apellido que quien en 1916 la ejecuta, con cuatro años de diferencia y sobre el mismo
inmueble. Que sean la misma persona es probable y no está acreditado, y se resuelve con el
expediente, no con el Boletín.
\end{cautela}

\textbf{Cuatro
eslabones, dos de ellos nuevos, y un solo dato de superficie en toda la serie.} Lo que falta ---quién
compró en la subasta de junio de 1916, si se realizó--- es exactamente lo que el expediente diría y
el Boletín no dice: ningún aviso de este departamento publica jamás el resultado de un remate.

Así se describía la propiedad rural del departamento en 1910: por lomas, cuchillas,
mojones antiguos y el nombre de quien poseía al lado. Sin coordenadas, sin superficie
declarada en el edicto, sin plano publicado. Ése es el material del que descienden las
matrículas que este libro persigue, y la razón de fondo por la cual establecer si una
proviene de otra exige un estudio de títulos y no una consulta.

\section{Una finca con cadena completa}

Hay un caso —uno solo— en el que este libro puede seguir un inmueble de punta a punta.

En enero de 1921 se ordena el deslinde de las fincas Cerro de Buena Vista y Sausal o Curuzú, en el partido de Vaqueros, de propiedad de Francisco S. Urquiza\index{Urquiza, Francisco S.}; en 1924 el mismo Urquiza obtiene concesión para regarlas con sobrantes del Wierna; y en 2014 la Secretaría de Ambiente convoca a audiencia por la Urbanización Cerro de Buena Vista sobre las matrículas 1713 y 1905. El capítulo \ref{cap:loteo} sigue esa secuencia.

Deslinde en 1921, agua en 1924, urbanización en 2014. Es la secuencia completa —título,
agua, loteo— documentada sobre un mismo inmueble a lo largo de noventa y tres años. Este
libro no afirma que las tres recaigan sobre idéntica superficie: un deslinde de hace un
siglo pudo subdividirse muchas veces. Afirma que el nombre y el lugar coinciden, que el
expediente de 1921 existe y está identificado, y que ése es el punto de partida natural
para reconstruir la cadena. Para los loteos del capítulo \ref{cap:loteo}, en cambio, ese punto de
partida todavía no aparece.

\section{Wierna: el recorrido completo de un nombre}

En octubre de 1908, el Boletín publica una sentencia sobre unos terrenos denominados «La
Finca». Se discute el título de don Bautista Wierna, y el fallo menciona «el callejón que
girando de Norte a Sud conduce al Molino de Wierna y después al Rosario de Lerma». La
demandada gana por prescripción treintañal, y el tribunal razona que «no figura en autos el
título de compra»: el reivindicante no pudo adquirir más derechos que los que tenía su
vendedor.

En 1921, Wierna ya no es un apellido sino un circuito electoral: figura entre los parajes
de una mesa receptora del departamento, junto a El Acheral, El Charcal, Las Garrapatas,
Los Porongos, Mesadas, Potrero de Valencia, Quesería, Santa Rufina y Yacones.

Hoy es un río, una finca con tres matrículas sobre las que se determinó la línea de ribera
en 2015, y el lugar de dos canteras de áridos concesionadas sobre tierra fiscal.

Apellido, paraje, curso de agua, unidad catastral, concesión minera. Ese recorrido es el
que siguen casi todos los topónimos de este departamento, y explica en una línea por qué
reconstruir su historia dominial exige leer un siglo de actos y no consultar un registro.

Vale la pena detenerse en cómo se resolvió aquel pleito de 1908. Se resolvió porque el
título no aparecía. La opacidad dominial del corredor no es un efecto del mercado
inmobiliario reciente: en el primer año del Boletín Oficial ya había un juicio que se
decidía por ausencia de instrumento traslativo de dominio, sobre la misma cuenca donde hoy
no puede establecerse si una matrícula proviene de otra.

\begin{cautela}
Hay que leer ese pleito con una reserva que la segunda lectura de 1908 hizo evidente. \textbf{La
sentencia no dice en qué departamento están los terrenos}, y el único rumbo que da es el de un
callejón que lleva al Molino de Wierna «y después al Rosario de Lerma», que queda al sudoeste de
la capital, mientras el río Wierna corre al norte. Puede tratarse de otra propiedad de un Wierna, fuera
de este departamento. Que el pleito sea «sobre la misma cuenca» es, entonces, una posibilidad y
no un dato; lo que se sostiene sin ella es que en 1908 un juicio de tierra se decidía por la
falta del título.
\end{cautela}

\section{El deslinde de 1932}

El documento más completo de este relevamiento entre los que nombran a la vez las fincas, sus
dueños y sus linderos es un edicto de deslinde de mayo de 1932, y conviene leerlo entero.

\begin{ficha}{Edicto de deslinde --- B.O.\ Nº 1.429, 27 de mayo de 1932}
Ángel R.\ Bascari, con poder y títulos de don \textbf{Silvano Y.\ Murúa} y de su esposa doña
Fanny Ovejero Lacroix de Murúa, solicita deslinde, mensura y amojonamiento de cinco fincas
del departamento La Caldera: \textbf{«Las Garrapatas», «Cañada Seca o Peñones», «El Durazno o
Peñones», «Los Matos» y «Severino»}.

Perito propuesto: Hermann Pfister\index{Pfister, Hermann}. Juzgado de 1ª Instancia y 1ª Nominación en lo Civil, Dr.\
Néstor Cornejo Isasmendi. Edictos por treinta días.

\textbf{Linderos nombrados en el edicto:} el puesto La Despensa, de los señores Adolfo,
Alberto y Abel Murúa; \textbf{la finca y la estancia Los Porongos, de don Daniel Linares\index{Linares, Daniel}};
el puesto Severino; la finca Peñones, de Silvano Y.\ Murúa; Peñones, de Francisco Sosaya ---«Sosauya», así en el edicto--- y
\textbf{herederos de Liborio Guerrero\index{Guerrero, Liborio}}; los herederos Campos; la sucesión de Silvano Murúa\index{Murúa, Silvano},
antes de Salvador Murúa; Ceibal o Algarrobal, de la sucesión del doctor José María Solá; El
Acheral, de los herederos Toledo; El Nogalito, de Salvador Rosa; y El Paraíso.

\textbf{Accidentes citados como límites:} el camino de la Punta de Agua, la Cañada de la
Iglesia, la mesada del Panteón, La Laguna, el Morro del Sunchal y la encrucijada del
Acheral.
\end{ficha}

\subsection*{El Durazno es una finca de 1932}

La tercera de las cinco se llama \textbf{«El Durazno o Peñones»}.

Es el nombre del loteo cuyos arroyos se determinaron en 2014, cuya agua para cuatrocientos
lotes se tramitó en 2015, cuyo plano se aprobó en 2021 y cuyas viviendas se inundaron en
2018 y 2019. En 1932 era una de las cinco fincas cuyo deslinde pidieron Silvano Y.\ Murúa y su esposa ante un juez civil ---un año después, el remate de La Despensa ya nombra como vecinos a la sucesión y a los herederos de Silvano Murúa, el padre homónimo, muerto antes de agosto de 1927 (apéndice \ref{cap:personas})---, con
límites al norte con los herederos Campos y al oeste con Daniel Linares\index{Linares, Daniel}.

Este libro no afirma que la finca de 1932 y la matrícula del loteo sean la misma superficie:
noventa años y varias particiones median entre una y otra, y establecerlo exige el estudio
de títulos que el apéndice \ref{cap:pedidos} reclama. Afirma que \textbf{el nombre existe como unidad de
dominio desde 1932 y que su deslinde judicial está identificado}, con juzgado, perito y
fecha. Es, para el loteo más discutido de este libro, el punto de partida de la cadena que
hasta aquí faltaba.

Y hay algo más en esa lista de cinco. \textbf{Las Garrapatas} figura como paraje del circuito
electoral de 1921. \textbf{Peñones} es el partido policial de 1909 —«Peñones y Chalchanio»,
con Liborio Guerrero\index{Guerrero, Liborio} como comisario auxiliar— y la finca cuyo título perdido se repuso
judicialmente en 1926. Las cinco fincas de un solo deslinde recorren, entre ellas, la mitad
de los nombres que este libro persigue.

\subsection*{Y dos años después se partió}

En septiembre de 1934, \textbf{Alberto Adolfo Murúa y Abel Murúa} solicitan el deslinde del
\textbf{Puesto Despensa}, «parte integrante de la estancia La Despensa», de aproximadamente
\textbf{dos mil seiscientas hectáreas}, con el mismo perito que dos años antes: Hermann
Pfister\index{Pfister, Hermann}.

Los límites que declara el edicto (Boletín Oficial Nº 1.549, 14 de septiembre de 1934) muestran lo que había pasado en el medio. Al norte, «propiedad de don Francisco Sosaya en parte y en otra con el
puesto Cañada Seca o Peñones \textbf{adjudicado al heredero Silvano Murúa\index{Murúa, Silvano Ignacio}}». Al sur, «el
puesto Los Matos, \textbf{adjudicado a la heredera Eulalia Murúa de Cornejo}». Al poniente,
«el puesto Las Garrapatas y Severino, de esta misma sucesión».

Son cuatro de las cinco fincas del deslinde de 1932 y el puesto de La Despensa ---que en 1932 figuraba como lindero, a nombre de Adolfo, Alberto y Abel Murúa---, ahora repartidos entre herederos con
nombre y apellido. \textbf{La partición hereditaria de un dominio grande es, en este
departamento, el mecanismo por el cual la tierra empieza a dividirse}, y aquí se lo ve
funcionar en tiempo real: en 1932 Silvano Y.\ Murúa y su esposa piden el deslinde del conjunto; en 1934 cada
puesto ya tiene adjudicatario entre los herederos de la sucesión de su padre, de la que él mismo es heredero (apéndice \ref{cap:personas}).

Es el mismo mecanismo que, sesenta años después, la literatura académica identifica como
origen del mercado de suelo de Vaqueros cuando se resuelve el juicio sucesorio de los
descendientes de Serrey\index{Serrey, Carlos}. No es una analogía: es el mismo procedimiento, en la misma
provincia, sobre el mismo tipo de dominio.

\subsection*{Y en 1936 se ejecutó}

Dos años después de repartida, la finca principal de la sucesión salió a remate por deuda de
los propios herederos.

\begin{ficha}{Remates judiciales de la finca La Despensa --- 1936}
El martillero Santiago Esquiú vende, por disposición del Juez de Comercio y en autos
seguidos por \textbf{Salvador Rosa contra Abel, Adolfo y Alberto Murúa}, la finca \textbf{La
Despensa}, en el partido del Sauce, departamento La Caldera.

\textbf{Límites:} «Norte, Francisco Sosaya y \textbf{Silvano Murúa\index{Murúa, Silvano Ignacio}}; Sud, \textbf{Eulalia
Murúa de Cornejo}; Este, José María Solá; Oeste, \textbf{Sucesión Murúa}».

Hubo \textbf{dos subastas}. La primera, el 6 de febrero de 1936, con base de \$23.333,32. La
segunda, semanas después, con base de \$25.000 --- \textit{un aumento poco habitual entre
una subasta y la siguiente, que conviene verificar contra el original antes de citarlo}.
\end{ficha}

La secuencia completa de esta sucesión cabe en cuatro años y muestra el mecanismo entero.

En \textbf{1932}, Silvano Y.\ Murúa y su esposa mandan deslindar cinco fincas, entre ellas
El Durazno o Peñones. En \textbf{1934}, los herederos Alberto Adolfo y Abel Murúa ---así, sin coma y dos veces, en el edicto; el deslinde de 1932, el remate de 1933 y la ejecución de 1936 nombran a tres personas: Adolfo, Alberto y Abel--- piden el deslinde del Puesto
Despensa y los límites ya declaran a cada puesto «adjudicado» a un heredero con nombre y
apellido. En \textbf{1936}, tres de esos herederos son ejecutados por un acreedor y la finca
mayor sale a remate, con los otros herederos como linderos.

Deslinde, partición, ejecución. \textbf{Es el ciclo completo por el cual un dominio grande se
fragmenta}, y en este caso se lo puede seguir entero en el Boletín Oficial porque las tres
etapas dejaron aviso público.

Ese ciclo importa por lo que viene después. La literatura académica identifica la resolución
de un juicio sucesorio como el origen del mercado de suelo de Vaqueros en los años noventa;
este libro documenta, sesenta años antes y sobre las fincas del otro municipio, el mismo
procedimiento funcionando de principio a fin. \textbf{La partición hereditaria no es un
antecedente del mercado de suelo caldereño: es su forma más antigua.}

\subsection*{Cuando una finca cambia de nombre}

En junio de 1937 se publica el deslinde de un inmueble llamado \textbf{«La Helvecia», antes
«Calderilla»}, en el departamento La Caldera. Es el documento más denso del período y lo es
por cuatro razones distintas.

\begin{ficha}{Deslinde, mensura y amojonamiento --- B.O.\ Nº 1.691, aviso Nº 3612}
Solicita \textbf{Hermann Pfister\index{Pfister, Hermann}}. Inmueble: \textbf{«La Helvecia», antes «Calderilla»}.

\textbf{Linderos:} «Norte, propiedad de \textbf{Augusto Regis\index{Regis, Augusto}}; Sud, propiedad de don
\textbf{Genaro Molina\index{Molina, Genaro}}; Naciente, herederos de don \textbf{Pedro Valencia}; Poniente, el
\textbf{Río de La Caldera}.»

Juzgado de 1ª Instancia y 1ª Nominación en lo Civil, Dr.\ Guillermo F.\ de los Ríos. Perito:
agrimensor Jorge de Bancarel. Auto del 10 de abril de 1937, conforme los artículos 570 y 474
del Código de Procedimientos. Publicación por treinta días en dos diarios y \textbf{por una
vez} en el Boletín Oficial.
\end{ficha}

\textbf{Primero: el topónimo se abandona y el nombre propio lo reemplaza.} «Calderilla» es un
paraje del departamento desde 1909 --- partido policial con comisario propio --- y sigue
nombrando hoy a una localidad y a un centro vecinal. La finca que lo llevaba pasa a llamarse
La Helvecia. En los mismos boletines, Vialidad financia los caminos «Calderilla al Desmonte
por Gallinato» y «Calderilla a Güemes por Gallinato»: \textbf{el nombre que la propiedad
abandona sigue nombrando la traza vial}. El topónimo sobrevive al dominio que lo usaba, que
es exactamente lo contrario del recorrido de Wierna.

\textbf{Segundo: el lindero al poniente es el río.} No un vecino: el curso. Una finca cuyo
límite occidental es el río La Caldera, deslindada judicialmente en 1937, en el mismo
departamento donde ochenta y cuatro años después una sentencia obligará a determinar qué se
construyó sobre las márgenes.

\textbf{Tercero: el lindero al norte preside el municipio.} \textbf{Augusto Regis\index{Regis, Augusto}} aparece
como propietario colindante ya en 1909 ---el remate de la décima parte del Angosto de Arrieta
lo nombra, con Julio Mercado y Susana Aguilera, como lindero al este--- y otra vez en junio de 1937, y veinte días después un decreto lo nombra
Presidente de la Comisión Municipal del Distrito de La Caldera. Es el mismo Regis\index{Regis, Augusto} que figura
como jurado suplente en 1908, que integró la comisión de defensas contra el río en 1910, que
fue designado miembro del cuerpo municipal en 1922, que lo presidía en 1933 y a quien
Vialidad encomendó en ese año la limpieza de un camino. \textbf{Treinta años de presencia
continua, y ahora también como titular de tierra lindera a la finca que se mensura.}

Y \textbf{Genaro Molina\index{Molina, Genaro}}, el lindero al sur, es el nombre del comisario auxiliar del partido
de La Calderilla en el decreto de 1909.

\textbf{Cuarto: el perito de 1932 es ahora el solicitante.} Hermann Pfister\index{Pfister, Hermann} había sido
designado perito en el deslinde de las cinco fincas de los Murúa en 1932 y en el del
Puesto Despensa en 1934. En 1937 pide el deslinde en nombre propio o por poder. En este mismo
corpus aparece además como inspector auxiliar de minas. \textbf{El técnico que mide el
departamento es, cinco años después, parte en una mensura.}

\begin{cautela}
Como siempre, coincidencias de nombre y no identificaciones registrales. Pero conviene
observar la acumulación: en un solo edicto de 1937, dos de los cuatro linderos llevan apellidos que el archivo registra en cargos públicos del departamento ---el presidente municipal y el comisario de un partido---, otro lindero es el río, y el solicitante es el perito oficial de las mensuras.

Este libro no extrae de ahí ninguna irregularidad. Extrae la constatación que la tercera
tesis enuncia: \textbf{durante la primera mitad del siglo, el conjunto de personas que
administraban el departamento y el conjunto de las que eran dueñas de su tierra se
superponían casi por completo.} Y una consecuencia práctica: quien quiera reconstruir el
dominio de este territorio hará bien en empezar por las nóminas de funcionarios.
\end{cautela}

\subsection*{Y una finca con folio, asiento y libro}

En febrero de 1937 se remata otra finca del departamento, y el edicto trae un dato que
ninguno de los anteriores traía.

\begin{ficha}{Remate judicial --- B.O.\ Nº 1.670, edicto Nº 3415}
Por orden del Juez de Comercio, en autos «Embargo Preventivo \textbf{Zulema Terán vs.\ Ramón
I.\ Juárez}», el martillero José María Decavi rematará con base de \$2.333,33 la finca
\textbf{«Chalchanio»}, departamento La Caldera, «con extensión y linderos le asignan sus
títulos registrados \textbf{Folio 342, asiento 438, Libro B La Caldera}».
\end{ficha}

Dos cosas.

\textbf{La finca está inscripta, y el edicto lo cita.} Folio, asiento y libro: la anotación
registral completa. Es la primera vez en las ediciones leídas que un aviso sobre el departamento remite al
registro en lugar de describir la propiedad por lomas y mojones. Ocho años después de la
catastración de 1929, el sistema registral funciona y los edictos lo usan.

Eso importa para el argumento del capítulo \ref{cap:opacidad}. Cuando este libro no puede establecer si una
matrícula proviene de otra, no está describiendo un territorio sin registrar: está
describiendo un registro que existe, funciona desde hace un siglo y no es consultable sin
arancel ni presencia física.

\textbf{Y «Chalchanio» es otro partido de 1909.} Figuraba entonces como «Peñones y
Chalchanio», con Liborio Guerrero\index{Guerrero, Liborio} como comisario auxiliar. Con éste, la lista de partidos de
1909 que reaparecen en el archivo como fincas con dueño, deuda o ejecución llega a seis:
San Alejo, Los Porongos, Potrero de Gallinato, La Despensa, \textbf{Vaqueros} ---rematada en
1935 como finca «Vaqueros o Entre Ríos» en la ejecución hipotecaria del Banco Provincial--- y
el partido de Peñones y Chalchanio, cuyos dos nombres reaparecen como fincas distintas.

\begin{cautela}
El ejecutado se llama \textbf{Ramón I.\ Juárez}, y en las mismas ediciones el comisario de
policía de La Caldera se llama \textbf{Ramón Juárez\index{Juárez, Ramón}}. El corpus no permite afirmar que sean
la misma persona ni descartarlo: la inicial del segundo nombre aparece sólo en el edicto.

Se consigna con esa reserva porque, de tratarse del mismo, sería otro caso de lo que este
libro viene documentando --- el comisario del departamento como propietario de una de sus
fincas --- y porque se resuelve leyendo el folio 342 del Libro B.
\end{cautela}

\textbf{Y ese folio vuelve a aparecer ocho años después, con otro ejecutado y con la finca ya
partida.}

\begin{ficha}{Remate judicial Nº 875 --- B.O.\ Nº 2.313 a 2.326, 21 de junio a 7 de julio de 1945}
Martillero \textbf{Leoncio M.\ Rivas}, por disposición del Juzgado de Paz Letrado Nº 1 de la
Capital, Dr.\ Marcelo Quevedo Cornejo, en autos «\textit{Embargo Preventivo Strachan, Yáñez y
Cía.\ vs.\ Martín Herrera}». Subasta del 30 de julio de 1945, base \textbf{\$671,01}, importe
del crédito en ejecución.

Bien: «\textit{\textbf{CUARTA PARTE INDIVISA DE LA FINCA ``CHALCHANIO''}, con extensión de 490
hs.\ 9 as.\ 36 ctms2.}», limitando al Naciente con el alambrado que la divide de la finca
«Mosquera»; al Norte, desde el esquinero en el \textbf{Abra del Angosto}, pasando por la
\textbf{loma de cal}, hasta un mojón en el \textbf{Abra de Chalchanio}; al Sud y parte del
Poniente, con un filo que la divide de la finca \textbf{«La Despensa»}; y al Oeste y Noroeste,
«\textit{con la fracción que le correspondió a don \textbf{Ernesto Echenique}}».

«\textit{Título registrado a Folio 342, Asiento 438 del Libro B, de Inmuebles de La Caldera. La
hipoteca ha sido registrada a Folio 108, Asiento 180 del Libro A, de Gravámenes de La Caldera.}»
\end{ficha}

Es el mismo folio, el mismo asiento y el mismo libro que cita el edicto de 1937. Entre una
fecha y la otra el ejecutado cambió ---Juárez en 1937, Herrera en 1945--- y el objeto de la
ejecución se redujo de la finca a \textbf{una cuarta parte indivisa} de ella, con al menos un
comunero nombrado, Ernesto Echenique, y con La Despensa como lindero sur.

\textbf{Es el mismo ciclo que este capítulo describe para La Despensa}, corrido ocho años y
una finca. Y refuerza la razón para pedir ese folio: no hay que leerlo para resolver una
homonimia, sino para reconstruir al menos dos traslaciones de dominio y una partición sobre un mismo asiento.

\begin{cautela}
El aviso no nombra paraje, no da número de catastro y no dice en ningún momento que la finca
esté «en La Caldera»: la ubicación descansa enteramente en el libro de registro. Los mojones
que menciona ---Abra del Angosto, loma de cal, Abra de Chalchanio--- y las fincas colindantes
«Mosquera» y «La Despensa» no fueron contrastados contra ninguna cartografía.

Una anomalía que queda registrada y no explicada: \textbf{el aviso declara vigencia hasta el 27
de julio y deja de publicarse el 7}. Las veinte ediciones siguientes, hasta el 13 de agosto, no
lo contienen, y ningún aviso posterior del mismo martillero corresponde a este juicio. El
corpus no dice si la subasta se suspendió, se levantó el embargo o el aviso se discontinuó por
otra razón.
\end{cautela}

\subsection*{Un padrón de parajes, en 1940}

Hay un documento que este relevamiento no buscaba y que ordena casi todos los topónimos de
este libro. No es un plano, no es un catastro y no es un registro: es un decreto electoral.

\begin{ficha}{Decreto del 30 de enero de 1940 --- ubicación de las mesas receptoras de votos y sus circuitos}
Dictado por el Ministerio de Gobierno, Justicia e Instrucción Pública conforme los artículos
29 a 31 de la Ley 122 de Elecciones de la Provincia, para el comicio del \textbf{3 de marzo de
1940}. Republicado en cuatro ediciones, siempre en anexo sin foliar.

La Caldera es el \textbf{Colegio Electoral Nº 2}, con dos circuitos y tres mesas:

\textbf{Circuito 5, Mesa 1} --- Escuela Elemental Mixta, La Caldera. Nueve parajes:
\textit{Calderilla, Chalchanio, Helvecia, \textbf{Getsemaní}, Peñones, Potrero de Castilla,
Pueblo, San Alejo, Severino}.

\textbf{Circuito 5, Mesa 2} --- Iglesia Parroquial, La Caldera. Treinta y un parajes:
\textit{Abra de la Sierra, Acheral, El Angosto, Angosto de Arrieta, Angostura, El Bordo,
Calderilla, Campo Alegre, Cañada Ancha, Corral de Barranca, Churcal, Los Chusos, Despensa,
Despensilla, Higueritas, \textbf{Potrero de Gallinato}, \textbf{Las Lagunas}, Los Matos,
Mesada, Monte de las Garrapatas, \textbf{Mosquera}, Nogalito, \textbf{Potrero de Valencia},
\textbf{Los Porongos}, Quesería, Santa Gertrudis, \textbf{Santa Rufina}, Los Sauces, La Sierra,
Wierna, Yacones}.

\textbf{Circuito 6, Mesa 1} --- Oficina del Registro Civil, Vaqueros. Siete parajes:
\textit{\textbf{Entre Ríos}, Estación F.C.C.N.A., Mojotoro, \textbf{Lesser}, El Túnel,
Vaqueros, \textbf{Villa Urquiza}\index{Urquiza, Francisco S.}}.
\end{ficha}

\textbf{Cuarenta y siete entradas y cuarenta y seis nombres distintos ---Calderilla figura en dos mesas---, y casi todos los que este libro persigue están ahí.} Getsemaní
---la matrícula 4366, a la que el edicto de agua de 2015 vincula con el loteo El Durazno, aprobado sobre la matrícula 6.558 (capítulo \ref{cap:ribera})--- figura como paraje
poblado y con mesa asignada en 1940. Las Lagunas es la finca cuya posesión treintañal se
tramitará cinco años después. Chalchanio es la finca que el remate de 1945 volverá a llevar
a subasta sobre el mismo folio, y Mosquera y Despensa son los dos linderos que ese remate
nombra. Potrero de Valencia es el lindero este del terreno cuya posesión treintañal promueve
Pfister\index{Pfister, Hermann} en 1945 ---«propiedad de la sucesión de Micaela de Valencia», según el edicto---, y
los herederos de Pedro Valencia eran ya, en 1937, el lindero naciente de La Helvecia. Potrero de Gallinato, Los Porongos, San Alejo,
Santa Rufina y Campo Alegre son partidos de 1909 ---el último, desde 1911---, fincas rematadas o fracciones expropiadas
en 1972. Y en el circuito de Vaqueros están Entre Ríos ---el otro nombre de la finca Vaqueros
rematada en 1935--- y una \textbf{Villa Urquiza}.

Lo que este decreto aporta no es un dato sino \textbf{una escala}: en 1940 el Estado sabía, y
publicaba, en qué parajes vivía la gente del departamento y a qué mesa le correspondía votar.
Es la lista más completa de unidades territoriales pobladas que este relevamiento encontró
entre el padrón de partidos de 1909 y la cartografía catastral de 2026, a treinta y un años del primero y a ochenta y seis
de la segunda.

\begin{cautela}
Un paraje electoral no es una finca ni una unidad catastral: nombra dónde vive gente que vota,
no dónde termina un dominio. Que «Getsemaní» sea paraje en 1940 no dice quién era su dueño ni
qué superficie tenía, y este libro no lo infiere. La correspondencia entre estos cuarenta y
seis nombres y las matrículas actuales es, exactamente, uno de los pedidos que el apéndice
\ref{cap:pedidos} dirige a la Dirección General de Inmuebles ---y ahora puede formularse con
una lista cerrada en lugar de con trece partidos.
\end{cautela}

\subsection*{Y el registro no reemplazó a los mojones}

El edicto de Chalchanio de 1937 es el primero de las ediciones leídas que identifica un inmueble por folio,
asiento y libro en lugar de describirlo por lomas y mojones. Sería cómodo leer eso como un
umbral: antes se describía, después se registraba.

\textbf{No es así, y hay un edicto de 1945 que lo desmiente.}

\begin{ficha}{Edicto Nº 460 --- posesión treintañal de la finca «Las Lagunas» --- 23 ediciones, enero a marzo de 1945}
Promovido por el procurador Diógenes R.\ Torres en representación de don \textbf{Gregorio
Maidana}, ante el Juzgado de Primera Instancia, Segunda Nominación en lo Civil, Dr.\ Roberto
San Millán. Auto del 30 de diciembre de 1944.

El inmueble se describe así, «\textit{según el título de adquisición de don Domingo Bejarano a
don Marcos Arroyo}»: «\textit{Desde el límite Norte, que divide la propiedad del otorgante
Arroyo con la del señor José Manuel Fernández hasta dar con un arroyo que baja de las cumbres
altas de Naciente a Poniente, denominado de ``Los Comederos'' o de ``El Barbecho'', y de
Naciente a Poniente, desde las cumbres altas de las serranías de La Caldera hasta dar con el
\textbf{Río de Los Yacones}, enfrentando con el arroyo de ``Las Carretas''}».

El auto ordena oficiar «\textit{al Sr.\ Juez de Paz de La Caldera}» para recibir la información
ofrecida, y preguntar a las oficinas «\textit{si el inmueble de referencia cuya posesión se
pretende acreditar, afecta o no propiedad fiscal o municipal}».
\end{ficha}

Ocho años después del edicto que citaba folio y asiento, una finca entera del departamento se
identifica ante la justicia \textbf{por vecinos y por cursos de agua}: sin superficie, sin
número de catastro, sin partida. Lo único cuantificado en el aviso es el importe de la
publicación, sesenta y cinco pesos.

De los cuatro nombres del título ---Maidana, Bejarano, Arroyo, Fernández--- la cadena de
adquisición sólo cubre a dos, y \textbf{Maidana no está en ella}: eso es exactamente lo que la
posesión treintañal viene a suplir. La misma vía que el Estado usó en 1943 para regularizar el
predio de la Comisaría, y que Hermann Pfister\index{Pfister, Hermann} usará seis meses después
en La Calderilla.

\textbf{Las dos formas de describir conviven}, y conviven con la tercera que el capítulo
\ref{cap:ribera} encuentra setenta años más tarde: pares de coordenadas Posgar que unas
resoluciones publican y otras remiten a la sede.

\subsection*{Los dueños eran los que gobernaban}

La lista de linderos merece leerse contra el capítulo \ref{cap:siglo}, porque los nombres coinciden de una
manera que no admite ser tratada como casualidad.

\textbf{Daniel Linares\index{Linares, Daniel}} figura como propietario de la finca y de la estancia Los Porongos.
Es el mismo Daniel Linares\index{Linares, Daniel} que en enero de 1910 integró la comisión de siete vecinos
designada para proyectar y ejecutar las defensas del pueblo contra las crecientes del río;
que la presidió; que en 1911 rindió las cuentas del déficit; y que en abril de 1922 fue
designado por el Interventor Nacional miembro de la Comisión Municipal de La Caldera.

\textbf{Silvano Ignacio Murúa\index{Murúa, Silvano Ignacio}}, que pide con su esposa el deslinde de las cinco fincas, no es el único
Silvano Murúa del expediente: el mismo edicto nombra como lindero a «la sucesión de Silvano
Murúa\index{Murúa, Silvano}»: su padre, nacido en 1855, muerto antes de agosto de 1927, del que el
titular es heredero. El que renunció a la Comisión Municipal en 1909 ---«Silvano J.\ Murúa» en el
decreto, reemplazado por Abelardo Figueroa--- y que integró la comisión de defensas de 1910 es el
hijo, entonces de veintiséis años: había nacido en 1882 y se había casado dos meses antes de la
renuncia (apéndice \ref{cap:personas}).

\textbf{Liborio Guerrero\index{Guerrero, Liborio}}, cuyos herederos aparecen como linderos, fue comisario auxiliar del
partido de Peñones y Chalchanio según el decreto de 1909.

\begin{cautela}
Corresponde el mismo cuidado de siempre. Se trata de coincidencias de nombre y apellido en
documentos separados por veinte y treinta años, no de identificaciones registrales; en un
departamento de pocos miles de habitantes los apellidos se repiten, y una sucesión
puede llevar el nombre del causante durante décadas.

Pero la coincidencia no es de un nombre: es de tres, y los tres ocupan en 1909, 1910 y 1922
los cargos que el capítulo \ref{cap:siglo} documenta —comisión municipal, comisaría de partido, comisión
de defensas del río—, y aparecen en 1932 como propietarios o causantes de las fincas
principales del departamento.

Lo que este documento permite afirmar, entonces, es lo que el capítulo \ref{cap:politica} sostiene por otra
vía cuando encuentra a un propietario de finca convertido en senador por el departamento:
\textbf{en este territorio, durante la primera mitad del siglo, quienes administraban el
municipio, vigilaban los partidos y proyectaban las obras contra el río eran los dueños de
la tierra sobre la que se hacían}. No es una imputación --- en 1932 no había otra forma de
gobernar un lugar así --- sino la descripción de una estructura, y es el antecedente más
antiguo de lo que la tercera tesis de este libro llama un dispositivo anterior al mercado.
\end{cautela}

\section{Un dato menor que confirma quiénes estaban}

Antes de llegar a la sucesión que abrió el mercado de suelo de Vaqueros, conviene registrar
una pieza lateral y elocuente. En 1934, la Sala en lo Penal de la Corte de Justicia resolvió
una causa por hurtos cometidos, según la carátula, \textbf{en las casas de Serrey\index{Serrey, Carlos} y de
Linares}, en el departamento de La Caldera.

No aporta nada sobre el dominio y este libro no lo usa para eso. Aporta otra cosa: que en
1934 las dos familias que este capítulo y el \ref{cap:siglo} vienen siguiendo --- la que casi sesenta años después abriría por partición hereditaria el mercado de parcelas de Vaqueros, y la del
presidente de la comisión de defensas del río de 1910 --- tenían casa en el departamento y
eran identificadas por su apellido en la carátula de un expediente penal.

\subsection*{La Curia y Serrey, linderos en el pueblo}

En noviembre de 1939 la \textbf{Curia Eclesiástica de Salta} promueve juicio de posesión
treintañal sobre un terreno del pueblo de La Caldera, ante el juez civil Carlos Cornejo
Costas. El edicto declara medidas irregulares --- ciento veintiséis metros al norte,
treinta y tres al sur, con un martillo de diez metros al sudeste --- y estos linderos:

«\textbf{Norte, Dr.\ Carlos Serrey\index{Serrey, Carlos}}; Sud, herederos Vidaurre y calle pública; Este, calle
pública; \textbf{Oeste, Dr.\ Carlos Serrey\index{Serrey, Carlos}}.»

Dos elementos de ese edicto conviene retener, porque los dos reaparecen ochenta años
después en municipios distintos del mismo departamento.

\textbf{La Iglesia adquiere por posesión treintañal.} No compra: acredita treinta años de
posesión pública y pacífica y pide que se le reconozca el dominio. Es el mismo instituto
que en 2015 usará la Provincia para regularizar el terreno de una escuela, y el mismo por el
cual en 1908 una demandada ganó un juicio de reivindicación porque el título del actor no
aparecía. En un territorio donde los títulos escasean, \textbf{la prescripción adquisitiva
no es un remedio excepcional sino una vía ordinaria de acceso al dominio.}

Y ochenta años más tarde, en Vaqueros, será también la Iglesia --- la Congregación de los
Hijos de la Inmaculada Concepción --- la que transfiera veintitrés hectáreas a un
fideicomiso para un desarrollo de doscientas veintisiete parcelas, según documenta el
capítulo \ref{cap:vaqueros}.

\textbf{Y el lindero, por dos costados, es Carlos Serrey\index{Serrey, Carlos}.} El mismo apellido cuya sucesión,
resuelta a principios de los años noventa, la literatura académica identifica como origen
del mercado de suelo de Vaqueros; el mismo que donó el terreno donde se levantó el edificio
municipal de ese municipio; y el mismo que llevaba Elena Serrey\index{Serrey, Elena} de González Bonorino, a
quien en 1986 se le expropió un inmueble de La Caldera para viviendas familiares ---su
parentesco con Carlos Serrey no consta (capítulo \ref{cap:expropiacion})---.

En 1939, Carlos Serrey\index{Serrey, Carlos} es dueño de los terrenos que rodean por dos lados a la Curia en el
casco del pueblo. \textbf{Los dos actores que en el presente aparecen convirtiendo suelo en
los dos municipios del departamento eran, hace ochenta y siete años, vecinos.}

\section{El padrón de 1944}

Hay un solo documento en todo este relevamiento que dice quiénes eran los propietarios del
departamento y cuánto valía lo de cada uno. Es de diciembre de 1944 y ocupa dos páginas del
Boletín Oficial.

El Jurado de Valuaciones publicó el \textbf{Revalúo General de la Provincia}, departamento por
departamento. El de La Caldera trae \textbf{cuarenta y siete partidas}, de la número 56 a la
184, con nombre y valuación. Los titulares son de todo tipo: personas, sucesiones
(«herederos de»), condominios («y otros») y una partida cuyos titulares eran menores de
edad, cuyos nombres este libro no reproduce por el criterio que fija el capítulo
\ref{cap:metodo}.
No hay decreto que lo introduzca ni subtotal que lo cierre: la
nómina entra directamente después de los avisos judiciales.

\textbf{Las trece mayores concentran el grueso del valor:}

\needspace{8\baselineskip}
{\sffamily\bfseries\footnotesize Revalúo General de la Provincia, departamento de La Caldera, diciembre de 1944\par}
\vspace{0.5em}
{\small
\setlength{\LTpre}{0.6em}\setlength{\LTpost}{1.2em}
\begin{longtable}{@{}p{1.4cm}p{7.4cm}p{3.0cm}@{}}
\toprule
\textbf{Partida} & \textbf{Titular} & \textbf{Valuación} \\
\midrule
\endhead
157 & \textbf{Urquiza\index{Urquiza, Francisco S.}, Francisco S.} & \$52.400 \\
104 & Rauchi, Celina Z.\ de & \$36.500 \\
80 & Sánchez, Juana Mercado de, y otros & \$26.400 \\
159 & Solá, Elvira Patrón Costa de & \$19.700 \\
82 & Bruzzo, José A.\ y Carlos M. & \$14.000 \\
170 & Ortiz, Elina Solá Patrón de & \$14.000 \\
152 & Yufra, Rosario & \$10.200 \\
96 & Ochoa, herederos de Concepción & \$8.900 \\
149 & Vivero, Sandalio & \$8.500 \\
169 & Berardi, Atilio & \$8.000 \\
88 & \textbf{Murúa, Silvano Ignacio} & \$7.600 \\
67 & Figueroa, Marcos B. & \$7.000 \\
105 & \textbf{Rejis} [\textit{sic}]\textbf{, Augusto} & \$6.400 \\
\midrule
\multicolumn{2}{@{}l}{\textit{Total de las 47 partidas publicadas}} & \textbf{\$353.100} \\
\bottomrule
\end{longtable}
}

Cinco cosas se leen ahí.

\textbf{Primera: Urquiza\index{Urquiza, Francisco S.} sigue siendo el mayor propietario de los que el revalúo publica.} Cincuenta y dos
mil cuatrocientos pesos, casi el quince por ciento de todo el valor publicado, y más del
cuarenta por ciento por encima del segundo. Veinte años después de la concesión del Wierna,
dieciséis después de su prisión preventiva y ocho después de perder el juicio de daños, su
posición patrimonial no se movió.

\textbf{Segunda: Murúa y Regis\index{Regis, Augusto} siguen ahí.} Silvano Ignacio Murúa, doce años después del
deslinde de las cinco fincas y ocho después del remate de La Despensa ---es la misma persona que el edicto de 1927 nombra \textbf{Silvano I.\ Murúa} y el de 1932, \textbf{Silvano Y.\ Murúa}: la inicial media cambia de un aviso a otro, pero el hombre es uno, el mismo que en 1924 y en 1932 pide deslindes junto con su esposa, María Fanny Ovejero de Murúa, y que en 1934 recibe como heredero el puesto Cañada Seca o Peñones---. Y «Rejis, Augusto»
--- así impreso --- que este libro identifica con el Augusto Regis\index{Regis, Augusto} que preside la comisión municipal y que el año anterior había dejado el cargo de comisario.

\textbf{Tercera: aparecen los apellidos de la elite provincial.} Elvira Patrón Costa de Solá y
Elina Solá Patrón de Ortiz suman treinta y tres mil setecientos pesos, casi el diez por
ciento del departamento. Son los apellidos que gobernaron Salta durante la década, y que
firman varios de los decretos que este libro transcribe.

\textbf{Y uno de esos apellidos ya era propietario veintisiete años antes.} Los dos deslindes de
la Calderilla publicados en 1917 los nombran: el del terreno de Juana Flores de Chaile linda al sud
con \textbf{Ildefonsa Burgos de Yugra} y al este con \textbf{la familia Yugra}, y el del Potrero
Gallinato, del mismo año, linda al oeste con \textbf{los herederos Yugra} (apéndice
\ref{cap:dominio}). No son titulares nombrados en un padrón fiscal: son linderos en un edicto
judicial, que es la forma más antigua en que este archivo registra una propiedad. \textbf{El revalúo
de 1944 no los encuentra: los vuelve a encontrar.}

\textbf{Cuarta: y aparecen los otros.} Mamani, Yugra, Apaza, Condori, Yavi, Yufra: apellidos de origen andino, con valuaciones que en su mayoría van de doscientos a cuatro mil pesos ---Rosario Yufra, con diez mil doscientos, es la excepción y figura entre las trece mayores---. \textbf{El padrón
registra las dos puntas del mismo territorio en la misma página.}

\begin{cautela}
\textbf{Y no aparecen en 1944: el archivo los tenía desde mucho antes, y con título.} El barrido
año por año del Boletín entre 1913 y 1932, que se hizo después de escribir este apartado, los
encuentra veintisiete años antes del revalúo y en las tres calidades que importan (capítulo
\ref{cap:siglo}):

\textbf{Como propietarios con deslinde judicial.} En 1917 se deslinda por el rumbo sud un terreno
en \textbf{La Calderilla} «\textit{de propiedad de doña \textbf{Juana Flores de Chaile}}», que
linda al sur con «\textit{terrenos de \textbf{Ildefonsa Burgos de Yugra}}», al este con
«\textit{propiedad de \textbf{familia Yugra}}» y al oeste con el río de la Calderilla (B.O.\ Nº
662). \textbf{Yugra es uno de los seis apellidos de este cuadro}, y en 1917 la familia tiene tierra
lindera con título suficiente para que un juez mande deslindar contra ella.

\textbf{Como causantes de sucesión ante el juez de paz del propio departamento.} El juicio
sucesorio de \textbf{don Delfín Chaile} se abre en octubre de 1918 por auto del Juez de Paz
propietario de La Caldera; los de \textbf{Simón Yufra} en 1924 y \textbf{Milagro Catacata} en 1927
se dictan y se firman en el pueblo.

\textbf{Y como autoridad.} El padrón de comisarios auxiliares de 1919 pone a \textbf{Delfín
Chaile} al frente de la Calderilla, a \textbf{Zacarías Maidana} en Potrero de Castilla y a
\textbf{Juan Cuchui} en Yacones; en 1929 el comisario de policía del departamento es
\textbf{Manuel Condori}, otro de los seis.

\textbf{Eso cambia lo que el cuadro de 1944 prueba.} No muestra la aparición de esas familias en
el padrón: muestra que en 1944 seguían ahí. \textbf{Y el descarte conviene decirlo también}: los
Apaza, Mamani y Aramayo que el barrido encuentra en 1919 y 1920 son de Guachipas, Carahuasi e
Iruya, y no de este departamento; el propio informe de barrido los identifica y los descarta.
Este libro no los cuenta.
\end{cautela}

\textbf{Y quinta: cuarenta y siete partidas para todo un departamento.} Si esas partidas fueran todas ---y la advertencia que sigue explica por qué no puede asegurarse---, un departamento de más de cien mil hectáreas estaría repartido entre menos de cincuenta titulares. Es la
estructura que el capítulo \ref{cap:tierra} encuentra setenta y cuatro años después en el Censo Agropecuario, cuando
una sola explotación concentra el cincuenta y uno por ciento de la superficie agropecuaria.
\textbf{La concentración no es un resultado del ciclo inmobiliario: es el punto de partida.}

\begin{cautela}
Dos advertencias sobre este cuadro.

La nómina publica \textbf{valuaciones fiscales, no superficies}. Un valor alto puede
corresponder a mucha tierra de poco valor o a poca de mucho, y el Boletín no distingue. Para
convertir estas cifras en hectáreas hace falta el padrón catastral, que el apéndice \ref{cap:pedidos} reclama.

Y las partidas no son correlativas: van de la 56 a la 184 con huecos. Eso puede significar que
la numeración es provincial y no departamental, o que se publicaron sólo las partidas
revaluadas. \textbf{Este libro no sabe cuál de las dos cosas es}, y por lo tanto no puede
afirmar que estos sean todos los propietarios del departamento en 1944.
\end{cautela}

\section{La sucesión Serrey}

La literatura académica sobre las transformaciones urbanas de Vaqueros sitúa la gran
expansión de loteos en la década del noventa e identifica, entre los factores de su origen,
uno de carácter estrictamente jurídico: la resolución de un largo juicio sucesorio de los
descendientes de Serrey\index{Serrey, Carlos}, que permitió regularizar a principios de esa década una parte
importante de los dominios y facilitó así la subdivisión y venta de nuevas parcelas.

Al mismo tiempo, el edificio de la Municipalidad de Vaqueros se levantó, bajo dirección de
Humberto Bini, en un terreno donado por la sucesión del doctor Carlos Serrey\index{Serrey, Carlos}. El municipio
había sido creado por ley provincial el 5 de agosto de 1970, y su primer intendente, Julio
Catalán Arellano, ejerció entre octubre de ese año y mayo de 1973.

\begin{cautela}
El mercado de suelo del municipio más dinámico del departamento no nace de una decisión de
política pública ni de una inversión externa: nace de un hecho de derecho sucesorio, la
partición de un dominio patricio históricamente indiviso. Antes de la sentencia, la tierra
existía como patrimonio familiar; después, como oferta.

Y el mismo linaje donó el terreno donde funciona el organismo que regula esa oferta. No
implica irregularidad alguna y no debe leerse como imputación. Implica que el municipio de
Vaqueros nace institucionalmente alojado en tierra donada por la familia cuya partición
hereditaria constituye el objeto principal de su función regulatoria posterior.

El apellido aparece además en los dos municipios y en posiciones opuestas. En 1986, la Ley
provincial 6354 declaró de utilidad pública y expropió un inmueble de Elena Serrey\index{Serrey, Elena} de
González Bonorino en la localidad de La Caldera —matrícula 1.303, 8.158,71 m²— con destino
a la construcción de viviendas familiares; su parentesco con Carlos Serrey\index{Serrey, Carlos} no consta. En Vaqueros, el linaje abre el mercado privado de parcelas por vía sucesoria y cede además, por dos expropiaciones de 1985 y 1986, tierra para viviendas y para la planta potabilizadora; en La Caldera, el mismo apellido cede tierra al Estado por expropiación para vivienda social (capítulo \ref{cap:expropiacion}).
\end{cautela}

\subsection*{Lo que este capítulo establece}

Las unidades territoriales que este libro persigue tienen más de un siglo. No son
fracciones creadas por el mercado inmobiliario: son las divisiones administrativas
históricas del departamento, que el fraccionamiento contemporáneo hereda, subdivide y
renombra.

El apéndice \ref{cap:dominio} ordena en una sola tabla las cincuenta y ocho operaciones de dominio que el archivo registra entre 1909 y 1945 ---remates, deslindes, particiones y posesiones treintañales---, con su bien, su causa, su martillero y su precio. Leída en vertical, muestra lo que este capítulo documenta caso por caso: que en 1932 y 1934 no había quien comprara tierra en este departamento al precio de tasación, y que la vía para regularizar un título era la misma para el Estado, para la Iglesia y para un particular.

Ésa es la tercera tesis del libro, y aquí queda establecida su base documental. Cuando el
capítulo \ref{cap:ribera} pregunte si una matrícula proviene de otra, estará preguntando por el último
eslabón de una cadena que arranca antes de que existiera el Boletín Oficial que la
publicaría.

